El presente proyecto de investigación tiene como propósito fundamental desarrollar un entorno experimental controlado para evaluar la percepción estudiantil en torno al desarrollo de UX y game feel mediante Inteligencia Artificial generativa. En cuanto a los aspectos a desarrollar, el alcance técnico y funcional contempla el diseño, programación e implementación de una demo técnica de un videojuego de gestión de recursos enfocado en la administración de una cafetería y la atención de pedidos. La arquitectura de software y el ciclo principal del juego (game loop) se construirán sobre el motor de código abierto Godot, integrando interfaces de usuario y sistemas de retroalimentación. El desarrollo de la lógica del sistema se ejecutará bajo el marco metodológico durante un periodo máximo de 6 meses. Este proceso dará como resultado dos versiones del juego: una versión base programada de forma tradicional (con ayuda de IA en caso de ser necesaria), y una versión asistida donde la refactorización y optimización de la lógica del código se realizarán mediante técnicas de vibe coding con la herramienta Claude, reflejando las adaptaciones propuestas por el modelo probabilístico en favor de la jugabilidad. Complementariamente, se incluye en el desarrollo la validación de un instrumento psicométrico en Google Forms basado en la escala de Likert y en la Player Experience Inventory (PXI) para medir cuantitativamente las variables del experimento ciego.

Por otra parte, se establecen exclusiones explícitas para delimitar los elementos que quedan fuera del alcance del proyecto, garantizando así la viabilidad temporal del trabajo para que este no supere un periodo de ejecución de seis meses. De este modo, quedan descartadas todas aquellas mecánicas avanzadas de progresión que extiendan innecesariamente el desarrollo, tales como sistemas de economía compleja, expansión de infraestructura virtual o árboles de habilidades. Asimismo, el software se limitará rigurosamente a una demo de corta duración, por lo que se excluye la producción de campañas extensas, narrativas profundas y el diseño de múltiples niveles comerciales. Finalmente, no se contempla la creación de sistemas de red, conectividad o modalidades multijugador, ejecutándose el juego exclusivamente de manera local y en modo de un solo jugador, al mismo tiempo que se delega el procesamiento físico y de renderizado a los componentes nativos de Godot, descartando el desarrollo de infraestructura de software base.

En última instancia, la ejecución del proyecto está sujeta a restricciones operativas y técnicos específicas, así como a supuestos clave que garantizan un marco metodológico realista. Como restricción primordial, se exige se mantenga la jugabilidad principal del juego (administración de café), haciendo que el modelo de IAG pueda tomar decisiones de diseño que sean comparables con las decisiones tomadas en el método tradicional. Adicionalmente, se establece la restricción de alcanzar una muestra mínima de sesenta participantes válidos de la Universidad del Valle de Guatemala para asegurar la significancia del ANOVA mixto 2x2. Bajo este contexto, se asumen como supuestos relevantes que la muestra posee una familiaridad básica con interfaces digitales para mitigar el sesgo en la curva de aprendizaje, y que las plataformas tecnológicas utilizadas (Claude, GitHub, Google Forms y Godot) mantendrán la estabilidad y disponibilidad de sus servicios gratuitos durante la vigencia de la investigación. 
